% Two stages, and what each one can honestly prove.
\begin{tikzpicture}[font=\sffamily\scriptsize, x=1cm, y=1cm]

\node[chlabel, anchor=west] at (-6.8,3.9) {two stages, split by what each one can honestly prove};

% ================================================================ stage one
\node[layerlabel] at (-6.8,3.4) {stage 1: every commit, no hardware, under four minutes};
\node[ext, text width=22mm] (commit) at (-5.6,2.5) {a commit};
\node[kern, text width=26mm] (sim) at (-2.4,2.5) {the simulated part, with its bus controller};
\node[hw, text width=22mm] (vcan) at (0.8,2.5) {a virtual interface};
\node[ext, text width=26mm] (host1) at (4.0,2.5) {the host stack, unchanged};
\foreach \a/\b in {commit/sim, sim/vcan, vcan/host1}{ \draw[flow] (\a) -- (\b); }
\node[chlabel, anchor=north] at (-2.4,1.85) {both frame formats, bridged};

\node[regcell, fill=usrcol, minimum width=54mm, text width=52mm, inner sep=3pt,
      anchor=north west, align=left] at (-6.8,1.3)
  {\textbf{can assert:} protocol correctness, state reachability, bounds
   checking, message mapping, error paths, the update sequence, and anything a
   fake can fail.};
\node[regcell, fill=red!18, minimum width=54mm, text width=52mm, inner sep=3pt,
      anchor=north west, align=left] at (0.2,1.3)
  {\textbf{cannot assert:} loop period jitter, following error, interrupt
   latency, flash timing, reaction time, \textbf{or anything with a microsecond
   in it}.};

% ================================================================ stage two
\node[layerlabel] at (-6.8,-0.6) {stage 2: once a night, on the bench, under twenty-five minutes};
\node[ext, text width=22mm] (pwr) at (-5.6,-1.5) {power cycle};
\node[fw, text width=26mm] (flash) at (-2.4,-1.5) {flash over the bus, chapter 19};
\node[hw, text width=22mm] (node) at (0.8,-1.5) {the board};
\node[ext, text width=26mm] (host2) at (4.0,-1.5) {the runner, and the master};
\foreach \a/\b in {pwr/flash, flash/node, node/host2}{ \draw[flow] (\a) -- (\b); }
\node[chlabel, anchor=north] at (0.8,-2.15) {no probe, no scope, no analyser};

\node[regcell, fill=hwcol, minimum width=114mm, text width=112mm, inner sep=3pt,
      anchor=north west, align=left] at (-6.8,-2.8)
  {\textbf{carries everything real silicon is needed for}: the six assertions,
   the thirty cycle trajectory set, two circular tests, and forty injected
   faults from a written catalogue. \textbf{The simulator provides no motion
   metrics at all}, so every number in this chapter is computed by the rig
   rather than read from a tool, in both stages.};

\node[umlnote, text width=114mm, anchor=north west] at (-6.8,-4.75)
  {Splitting the stages this way is what makes the first one fast enough to run
   on every commit and the second one honest enough to be worth running at all.
   A rig that asserts microsecond timing in simulation reports good news, and a
   rig that needs a bench for every commit is a rig people learn to skip.};
\end{tikzpicture}
